% Lean: EconLean/Bargaining/Nash.lean, EconLean/Cooperative/{Shapley,Core}.lean
\chapter{Bargaining and cooperative games}
\label{chap:cooperative}

Sources: Nash (1950); MWG §22.E and §18.AA; Shapley (1953); Bondareva (1963); Shapley (1967);
Osborne--Rubinstein, \emph{A Course in Game Theory}, Ch.~13--14.

\section{Nash bargaining}

\begin{definition}[Bargaining problem, Nash solution]
\label{def:bargaining}
\lean{Econ.BargProblem, Econ.nashProduct, Econ.IsNashSolution, Econ.affMap, Econ.continuous_affMap, Econ.convex_image_affMap, Econ.BargProblem.affine, Econ.BargProblem.IsSymmetric, Econ.BargSolution}
\leanok
A compact convex $S \subseteq \mathbb{R}^2$ with disagreement point $d \in S$ and some $s \in S$ with $s > d$. The
Nash solution maximizes $(s_1 - d_1)(s_2 - d_2)$ over $s \in S$, $s \ge d$.
\end{definition}

\begin{definition}[Nash's axioms]
\label{def:nash-axioms}
\lean{Econ.IsWeakPareto, Econ.IsSymmetricSol, Econ.IsAffineInvariant, Econ.IsIIA}
\leanok
\uses{def:bargaining}
Weak Pareto efficiency, symmetry, invariance to positive affine transformations of utilities, independence of
irrelevant alternatives.
\end{definition}

\begin{lemma}[Existence and uniqueness of the Nash solution]
\label{lem:nash-solution}
\lean{Econ.exists_nashSolution, Econ.IsNashSolution.nashProduct_pos, Econ.nashSolution_unique, Econ.nashSol, Econ.nashSol_spec}
\leanok
\uses{def:bargaining}
The Nash product attains its maximum on the compact individually rational part of $S$, strictly above $d$, and the
maximizer is unique.
\end{lemma}
\begin{proof}
\leanok
Weierstrass. For uniqueness, if $a = n - d$ and $b = m - d$ both maximize with value $p > 0$, the midpoint is
feasible with product $(a_1+b_1)(a_2+b_2)/4 \le p$, which forces $a_1 b_2 = a_2 b_1$ and then $a = b$.
\end{proof}

\begin{lemma}[Tangent line at the normalized solution]
\label{lem:sum-le-two}
\lean{Econ.sum_le_two_of_max}
\leanok
If $d = 0$ and $s_1 s_2$ is maximized over the individually rational part of a convex $S$ at $(1,1)$, then
$s_1 + s_2 \le 2$ on $S$.
\end{lemma}
\begin{proof}
\leanok
Otherwise a small step from $(1,1)$ toward such $s$ raises the product above $1$.
\end{proof}

\begin{theorem}[Nash's bargaining theorem]
\label{thm:nash-bargaining}
\lean{Econ.nash_bargaining_theorem, Econ.eq_nashSol_of_axioms, Econ.nashSol_axioms}
\leanok
\uses{def:nash-axioms, lem:nash-solution, lem:sum-le-two}
A solution satisfies the four axioms iff it is the Nash solution.
\end{theorem}
\begin{proof}
\leanok
Normalize affinely so that $d = 0$ and the Nash solution is $(1,1)$. By the lemma, $S$ lies in the symmetric set
$T = \{s_1 + s_2 \le 2, s \ge -M\}$, where symmetry and Pareto give $f(T) = (1,1)$; IIA gives $f(S) = (1,1)$; undo
the normalization. Conversely the Nash product scales by $a_1 a_2 > 0$ under affine maps and is symmetric, and its
maximizer over a larger set, if feasible in a smaller one, maximizes there.
\end{proof}

\section{The Shapley value}

\begin{definition}[TU game, marginal contributions, Shapley value]
\label{def:shapley}
\lean{Econ.TUGame, Econ.PlayerOrder, Econ.predSet, Econ.marginalContribution, Econ.shapley}
\leanok
A TU game assigns a worth $v(S)$ to every coalition. The Shapley value of $i$ is the average over all arrival
orders of $v(P_i \cup \{i\}) - v(P_i)$, where $P_i$ are the players who arrive before $i$.
\end{definition}

\begin{definition}[Shapley's axioms]
\label{def:shapley-axioms}
\lean{Econ.IsSymmetricPair, Econ.IsNullPlayer, Econ.IsShapleyAxioms}
\leanok
\uses{def:shapley}
Efficiency $\sum_i \varphi_i(v) = v(N) - v(\varnothing)$, symmetry, null player, additivity.
\end{definition}

\begin{theorem}[The Shapley value satisfies the axioms]
\label{thm:shapley-axioms}
\lean{Econ.sum_marginal, Econ.shapley_efficient, Econ.map_swap_of_iff, Econ.IsSymmetricPair.map_swap, Econ.shapley_symmetric, Econ.shapley_null, Econ.shapley_add, Econ.shapley_axioms}
\leanok
\uses{def:shapley-axioms}
\end{theorem}
\begin{proof}
\leanok
Along each order the marginal contributions telescope to $v(N) - v(\varnothing)$. Composing orders with the swap of
two symmetric players maps the marginal contributions of one to the other. Null players contribute nothing, and the
value is linear in $v$.
\end{proof}

\begin{lemma}[Unanimity decomposition]
\label{lem:unanimity}
\lean{Econ.unanimity, Econ.dividend, Econ.sum_alternating, Econ.game_eq_sum_unanimity}
\leanok
\uses{def:shapley}
Every game is $v = \sum_T c_T u_T$ with $u_T(S) = [T \subseteq S]$ and Harsanyi dividends
$c_T = \sum_{R \subseteq T} (-1)^{|T| - |R|} v(R)$.
\end{lemma}
\begin{proof}
\leanok
Möbius inversion on the subset lattice: $\sum_{R \subseteq T \subseteq S} (-1)^{|T| - |R|} = [R = S]$.
\end{proof}

\begin{theorem}[Shapley 1953: uniqueness]
\label{thm:shapley-unique}
\lean{Econ.IsShapleyAxioms.map_zero, Econ.IsShapleyAxioms.map_sum, Econ.IsShapleyAxioms.apply_unanimity, Econ.IsShapleyAxioms.unique, Econ.eq_shapley_of_axioms}
\leanok
\uses{thm:shapley-axioms, lem:unanimity}
The Shapley value is the only value satisfying the four axioms.
\end{theorem}
\begin{proof}
\leanok
On $c\,u_T$, players outside $T$ are null and those in $T$ symmetric, so efficiency gives each member $c/|T|$.
Additivity extends this to every game through the unanimity decomposition.
\end{proof}

\section{The core}

\begin{definition}[Core, balanced families, balanced games]
\label{def:core}
\lean{Econ.IsInCore, Econ.IsBalancedFamily, Econ.IsBalanced}
\leanok
\uses{def:shapley}
$x$ is in the core if $\sum_i x_i = v(N)$ and $\sum_{i \in S} x_i \ge v(S)$ for all $S$. Weights $\lambda_S \ge 0$
are balanced if $\sum_{S \ni i} \lambda_S = 1$ for each $i$; $v$ is balanced if $\sum_S \lambda_S v(S) \le v(N)$
for every balanced family.
\end{definition}

\begin{lemma}[Auxiliary identities]
\label{lem:core-aux}
\lean{Econ.sum_weight_mul_sum, Econ.singletonWeights, Econ.sum_singletonWeights_mem, Econ.sum_singletonWeights_mul, Econ.singletonWeights_nonneg, Econ.matrixPayoff_degenerate_left, Econ.matrixPayoff_degenerate_right}
\leanok
\uses{def:core}
$\sum_S \lambda_S \sum_{i \in S} x_i = \sum_i x_i \sum_{S \ni i} \lambda_S$; weights on singletons; matrix
payoffs against pure strategies.
\end{lemma}
\begin{proof}
\leanok
Exchange the order of summation.
\end{proof}

\begin{theorem}[Bondareva--Shapley]
\label{thm:bondareva-shapley}
\lean{Econ.isBalanced_of_mem_core, Econ.exists_mem_core_of_pos, Econ.bondareva_shapley}
\leanok
\uses{def:core, lem:core-aux}
The core is nonempty iff the game is balanced.
\end{theorem}
\begin{proof}
\leanok
($\Rightarrow$) $\sum_S \lambda_S v(S) \le \sum_S \lambda_S x(S) = \sum_i x_i = v(N)$.
($\Leftarrow$) Adding $K|S|$ shifts the core and preserves balancedness, so assume $v(\{i\}) > 0$, hence $v(N) > 0$.
In the zero-sum game where the maximizer picks a coalition $S$ and the minimizer a player $i$, with payoff
$v(S)/v(N) - [i \in S]$, let $(\mu, y)$ be a saddle point (minimax theorem). If the value were positive, then
$c_i = \sum_{S \ni i} \mu_S < c = \sum_S \mu_S v(S)/v(N)$ for all $i$, and $\lambda = \mu/c$ plus weights
$1 - c_i/c$ on singletons is balanced with $\sum \lambda_S v(S) > v(N)$. So the value is $\le 0$, and
$x = v(N)\,y$ is in the core.
\end{proof}

\section{Rubinstein bargaining}
\label{sec:rubinstein}

\begin{theorem}[Rubinstein 1982; Shaked--Sutton 1984]
\label{thm:rubinstein}
\lean{Econ.rubinsteinShare, Econ.one_sub_mul_pos, Econ.rubinstein_stationary, Econ.shaked_sutton_bounds, Econ.rubinsteinShare_mem}
\leanok
With discount factors $\delta_1, \delta_2 \in [0,1)$ the proposer shares $x^* = (1-\delta_2)/(1-\delta_1\delta_2)$ and
$y^* = (1-\delta_1)/(1-\delta_1\delta_2)$ are the unique solution of $x = 1 - \delta_2 y$, $y = 1 - \delta_1 x$. If the sup and inf
of SPE proposer payoffs satisfy $m_1 \ge 1 - \delta_2 M_2$, $M_1 \le 1 - \delta_2 m_2$ and symmetrically, they all equal these
shares. (Deriving the four inequalities from strategies is not formalized.)
\end{theorem}
\begin{proof}
\leanok
Chaining gives $M_1(1 - \delta_1\delta_2) \le 1 - \delta_2 \le m_1(1-\delta_1\delta_2)$, so $M_1 = m_1 = x^*$; likewise for player 2.
\end{proof}

\begin{theorem}[Comparative statics and frequent offers; Binmore--Rubinstein--Wolinsky 1986]
\label{thm:rubinstein-limit}
\lean{Econ.rubinstein_first_mover, Econ.rubinstein_patience, Econ.tendsto_one_sub_exp_div, Econ.rubinstein_limit, Econ.rubinstein_symmetric_limit}
\leanok
\uses{thm:rubinstein}
With a common $\delta$ the proposer gets $1/(1+\delta) > 1/2$; a player's share rises with her own patience and falls
with her opponent's. With $\delta_i = e^{-r_i\Delta}$ and $\Delta \to 0^+$, player 1's share tends to $r_2/(r_1+r_2)$, the
asymmetric Nash solution with bargaining powers proportional to $1/r_i$; with $r_1 = r_2$ it tends to $1/2$.
\end{theorem}
\begin{proof}
\leanok
$x^* = \frac{(1-e^{-r_2\Delta})/\Delta}{(1-e^{-(r_1+r_2)\Delta})/\Delta}$ and $(1 - e^{-a\Delta})/\Delta \to a$ (derivative of $t \mapsto 1 - e^{-at}$ at $0$).
\end{proof}

\section{The assignment game}
\label{sec:assignment-game}

\begin{definition}[Assignment game]
\label{def:assignment-game}
\lean{Econ.IsPairMatching, Econ.IsMatchingIn, Econ.matchingValue, Econ.isMatchingIn_empty, Econ.assignmentGame, Econ.le_assignmentGame, Econ.exists_optimal_matching, Econ.assignmentGame_singleton, Econ.assignmentGame_nonneg}
\leanok
Buyers $B$ and sellers $S$; a pair $(b,s)$ creates surplus $a_{bs}$. A coalition $T \subseteq B \sqcup S$ is worth the
largest total surplus of a matching among its members (possibly empty).
\end{definition}

\begin{theorem}[Shapley--Shubik 1971: the core]
\label{thm:assignment-core}
\lean{Econ.matching_sum_le, Econ.assignment_inCore_iff, Econ.assignment_weak_duality, Econ.assignLin, Econ.assignLin_convexOn, Econ.assignLin_perm_le, Econ.assignment_isBalanced, Econ.assignment_core_nonempty, Econ.shapley_shubik, Econ.assignment_strong_duality}
\leanok
\uses{def:assignment-game, thm:bondareva-shapley}
A payoff vector is in the core iff it is nonnegative, covers every pair ($u_b + p_s \ge a_{bs}$) and sums to the optimal
assignment value; every such covering vector pays at least that value. The game is balanced, so the core is
nonempty, and it is exactly the set of optimal solutions of the dual of the assignment LP.
\end{theorem}
\begin{proof}
\leanok
Core conditions on singletons and pairs give nonnegativity and covering; conversely the optimal matching of $T$
is paid for by its members. For balancedness, let $X_{bs} = \sum_T w_T [(b,s) \in \mu_T]$ for optimal matchings $\mu_T$;
it is doubly substochastic, and $\begin{psmallmatrix} I - \mathrm{diag}(X\mathbf 1) & X \\ X^\top & I - \mathrm{diag}(X^\top \mathbf 1)\end{psmallmatrix}$
is doubly stochastic on $B \sqcup S$. By Birkhoff--von Neumann it is a mixture of permutation matrices, and the
buyer--seller entries of each permutation form a matching of everyone, so $\sum_T w_T v(T) = \sum a_{bs} X_{bs} \le v(B \sqcup S)$.
Nonemptiness follows by Bondareva--Shapley (directly if there are no players).
\end{proof}

\begin{theorem}[Lattice of the core and competitive prices (Shapley--Shubik 1971; Roth--Sotomayor 1990, Ch.~8)]
\label{thm:assignment-lattice}
\lean{Econ.IsCorePayoff, Econ.isCorePayoff_iff, Econ.IsCorePayoff.weak, Econ.core_sup_inf, Econ.isCompact_core, Econ.buyer_optimal_core, Econ.seller_optimal_core, Econ.IsOptimalMatching, Econ.BMatched, Econ.SMatched, Econ.matching_sum_fst, Econ.matching_sum_snd, Econ.core_slack, Econ.IsCompEq, Econ.buyerSurplus, Econ.compEq_core, Econ.core_compEq, Econ.compEq_prices_iff, Econ.exists_min_compEq_prices, Econ.unit_assignment_value}
\leanok
\uses{thm:assignment-core}
Write core payoffs as buyer payoffs $u$ and prices $p$ (net of sellers' reservation values). If $(u,p)$ and $(u',p')$
are in the core, so are $(u \vee u', p \wedge p')$ and $(u \wedge u', p \vee p')$. There is a buyer-optimal core point
(best for every buyer, worst for every seller) and a seller-optimal one. At a core point every pair of an optimal
matching splits exactly its surplus and unmatched agents get $0$. A competitive equilibrium $(p, M)$ (each matched
buyer gets a seller in her demand set, unmatched buyers demand nothing at positive surplus, unsold objects have price
$0$) has $M$ optimal and $(\text{surpluses}, p)$ in the core; conversely core prices with any optimal matching are
competitive. Hence competitive prices are exactly core prices, and the minimal competitive prices are those of the
buyer-optimal core point. Witness: one buyer and one seller with surplus $1$.
\end{theorem}
\begin{proof}
\leanok
Both lattice points cover every pair (compare the larger buyer payoff with the matching price) and their totals add
to twice the optimal value, while each is at least that value by weak duality. The core is compact, so each buyer's
best core payoff is attained; joining these finitely many points gives the buyer-optimal point, and its prices are
minimal since a lower price would make the meet's total fall below the optimal value. Slackness: the gap
$\sum_{(b,s)\in M}(u_b + p_s - a_{bs}) + \sum_{\text{unmatched}} \text{payoffs}$ is a sum of nonnegative terms equal to $0$.
For a competitive equilibrium the buyers' surpluses cover every pair and the total equals the value of $M$, which by
weak duality is optimal.
\end{proof}

\section{The alternating-offers game}
\label{sec:rubinstein-game}

\begin{definition}[Alternating-offers game, SPE]
\label{def:barg-game}
\lean{Econ.BargProfile, Econ.BargProfile.path, Econ.BargProfile.Agrees, Econ.BargProfile.payoff, Econ.BargProfile.respPayoff, Econ.BargProfile.path_length, Econ.BargProfile.path_succ', Econ.BargProfile.payoff_eq_resp, Econ.BargProfile.payoff_congr, Econ.BargProfile.payoff_le_of, Econ.BargProfile.payoff_nonneg, Econ.BargProfile.payoff_add_le, Econ.pay₁, Econ.pay₂, Econ.IsDev₁, Econ.IsDev₂, Econ.IsBargSPE}
\leanok
A history is the list of rejected offers (player 1's share); player 1 proposes after histories of even length,
player 2 after odd ones. A pure strategy profile gives an offer in $[0,1]$ after every history and a response to every
offer. Agreement on $x$ after $t$ further periods gives $(\delta_1^t x, \delta_2^t (1-x))$ from the current node; perpetual
disagreement gives $0$. A profile is a subgame perfect equilibrium if at every proposal node and every response node
no player gains from changing her own choices.
\end{definition}

\begin{theorem}[Rubinstein 1982]
\label{thm:rubinstein-game}
\lean{Econ.stationaryProfile, Econ.stationary_eqs, Econ.stationary_pay₁, Econ.stationary_pay₂, Econ.dev₁_pay_le, Econ.dev₂_pay_le, Econ.stationary_isBargSPE, Econ.BargProfile.setAcc, Econ.BargProfile.setProp, Econ.BargProfile.setAcc_payoff_next, Econ.BargProfile.setAcc_acc, Econ.BargProfile.setProp_payoff, Econ.setAcc_isDev₁, Econ.setAcc_isDev₂, Econ.setProp_isDev₁, Econ.setProp_isDev₂, Econ.spe_acc₁, Econ.spe_acc₂, Econ.spe_acc₁_le, Econ.spe_acc₂_le, Econ.spe_prop₁, Econ.spe_prop₂, Econ.spePay₁, Econ.spePay₂, Econ.pay_bounds, Econ.spePay_props, Econ.spePay_eq, Econ.bargSPE_eq_stationary, Econ.bargSPE_unique}
\leanok
\uses{def:barg-game, thm:rubinstein}
The stationary profile (offer $x^*$ resp.\ keep $y^*$; accept iff one's share is at least $\delta_1 x^*$ resp.\ $\delta_2 y^*$) is
a subgame perfect equilibrium, and every pure-strategy SPE coincides with it after every history. In particular
player 1 offers $x^*$ at once and player 2 accepts.
\end{theorem}
\begin{proof}
\leanok
Against a stationary opponent every agreement gives the deviator at most the stationary share, so no deviation pays.
For uniqueness let $M_i, m_i$ be the sup and inf of proposer payoffs over all SPEs and subgames. A responder
must accept any offer beating her discounted continuation and may accept only offers at least as good; a proposer gets
at least any accepted offer. These give the Shaked--Sutton inequalities, so all four bounds equal the stationary
shares; a delayed agreement would give the proposer at most $\delta_i^2$ times her share, so agreement is immediate on the
stationary offer, and responses are pinned down by the strict comparisons plus the indifferent case.
\end{proof}
